\subsection{Summary of contributions}
We extended a computational sCPG model --- a two-legged half-centre
network in the Hodgkin--Huxley lineage
running from Rybak 2006 through Shevtsova 2015 and Danner 2017 to
Zhang 2022~\cite{rybak2006,shevtsova2015,danner2017,zhang2022} --- with four
mechanisms absent from that lineage: (i)
STDP that lets the network tune its own
cutaneous and descending synapses onto the rhythm generators rather than have
them hand-set; (ii) a closed loop from muscle-force feedback (Ia) into the
reciprocal-inhibition core; (iii) motoneuron pools driving explicit muscle
proxies with force, length and fatigue dynamics; and (iv) closed-loop stance
termination, in which each leg's stance ends when its own extensor force falls
rather than on a clock. Using this model we asked how a spinal
CPG's own synaptic weights should be organised to produce a stable walking
rhythm, and how that self-tuned rhythm depends on walking speed, on partial
unloading and on how fast the synapses can learn --- the variables manipulated
in rodent and human studies of SCI rehabilitation.

Three findings anchor the paper. First, the three plastic synapses --- the
cutaneous CUT$\to$RG-E projection and the brainstem BS$\to$RG-E and
BS$\to$RG-F projections --- converge to the same strength regardless of where
they start (from silent to $16$\,pA; \autoref{sec:results:robustness}), of
which leg they are in, and of walking speed or partial unloading
(\autoref{sec:results:weights}); the network organises its own gains rather
than requiring them to be prescribed by the modeller. Second, learning is what
closes the sensorimotor loop: once the cutaneous synapse has learned, stance is
ended by the leg's own force in every locomotion mode, whereas without learning
the network falls back on its safety timer (\autoref{sec:results:comparison}).
Third, the worst gait is produced not by the absence of learning but by a
partially learned synapse, strong enough to take over the timing of stance but
too weak to organise it (\autoref{sec:results:comparison}).

\subsection{Relation to prior computational and experimental sCPG work}
Zhang et al.\ (2022)~\cite{zhang2022} name two limitations of their own
four-rhythm-generator computational sCPG model explicitly: it has no
biomechanics or sensory feedback, and it has no motoneurons, treating
rhythm-generator output as a direct stand-in for motor output. Every
synaptic weight in their model, like its Rybak/Shevtsova/Danner
predecessors, is hand-tuned. Our model addresses both limitations for a
single hindlimb pair: the Ia force/length feedback loop supplies the
missing sensory closed loop, and explicit motoneuron pools driving a
saturating muscle model supply the missing motor periphery. On top of this
we add a mechanism none of these models have --- activity-dependent
plasticity on the afferent synapses --- so that the descending and
cutaneous gains are discovered by the network rather than fixed by the
experimenter.

The trade-off is architectural detail. Zhang et al.\ resolve genetically
identified commissural and long-propriospinal interneuron classes (V0V,
V0D, V1, V2a, V3, Shox2) across four limbs and reproduce V3-silencing
experiments quantitatively; our model collapses the commissural pathway to
two hand-set inhibitory projections between homologous half-centres and
covers only one limb pair. We also do not reproduce gait-transition
phenomena (walk $\to$ trot $\to$ gallop $\to$ bound): our stance timing is
closed-loop, but walking speed is set by per-mode fatigue and timeout constants
rather than emerging from a single tonic-drive parameter.
The two efforts are complementary rather than competing: theirs
demonstrates architectural realism at the cost of a closed loop; ours
demonstrates a closed, self-organising loop at the cost of architectural
detail. A natural synthesis --- hand-curated interlimb architecture wrapped
in a plastic, closed sensorimotor loop --- is a direct extension of both.

Our loading manipulation is drawn from a second, experimental lineage: the
rat and human spinal-cord-injury work of Courtine, Lavrov, Edgerton,
Gerasimenko and colleagues, in which stepping is studied across a graded
series of load-bearing conditions, from fully weight-supported walking to
hindlimb suspension~\cite{lavrov2008,edgerton2008,courtine2009,harkema2011}.
That literature establishes, largely from kinematic and electromyographic
readouts, that load-related afferent input shapes spinal stepping. Our model
offers a mechanistic, spiking-network account of one part of that picture: the
paw-contact signal both triggers the stance/swing transitions and trains the
synapse that makes those transitions reliable, so partial unloading weakens
the gait and slows its acquisition (\autoref{sec:results:weights},
\autoref{sec:results:comparison}). In that sense the model is positioned to
connect the computational sCPG lineage above with this experimental
SCI-rehabilitation lineage, and we return to that connection in the Clinical
implications below.

\subsection{A partially learned synapse is worse than none}
One finding deserves a mechanistic reading before its clinical one. The gait
is worst at $\lambda=10^{-5}$, where the cutaneous synapse has grown to only
$\approx7$--$14$\,pA by the end of the run, not at $\lambda=10^{-6}$, where it
has not grown at all (\autoref{sec:results:comparison}). With no learning the
extensor force rarely reaches the stance-termination threshold, so stance
bouts run to the timeout and the timeouts impose a regular, clock-like
alternation. With converged learning the cutaneous drive produces a strong,
well-timed extensor burst and the force threshold ends each stance reliably.
In between, the drive is strong enough that the force threshold, rather than
the timer, decides when stance ends, but too weak to make that decision
consistent from step to step, so both extensor/flexor and left--right
coordination degrade. The same logic applies along the loading axis: at
$\lambda=10^{-5}$ toe stepping, with half the cutaneous drive, has not even
reached that intermediate regime and still runs on the timer. In both cases
\emph{some} correctly organised plastic drive is better than none, but a
partially organised one can be worse than either extreme.

\subsection{Clinical implications}
The present computational approach yields several insights for the
clinical rehabilitation of individuals with SCI. Apart from recapitulating
locomotor behaviour, the proposed computational framework provides a
mechanistic account of how sensorimotor interactions and activity-dependent
plasticity might enable functional recovery. First, sensory inputs are
critical for facilitating locomotor rehabilitation, not as a replacement
for lost supraspinal descending input, but as a catalyst for plastic
adaptation in the spinal cord sensorimotor networks.

Our simulation results show that, when appropriately controlled through
load and/or activity, such sensorimotor input drives spinal adaptation and
subsequently restores coordinated locomotor outputs
(\autoref{sec:results:weights}).

Weight-bearing conditions drove the cutaneous synapse to its plateau
quickly and stabilised locomotor output, whereas reduced support slowed that
learning and led to a less coordinated gait
(\autoref{sec:results:weights}, \autoref{sec:results:comparison}). These findings parallel clinical
observations that repetitive task-specific locomotor rehabilitation ---
body-weight-support treadmill walking, robot- and manually-assisted
stepping --- restores stepping in individuals with motor-incomplete SCI by
engaging residual plasticity in the spinal sensorimotor circuitry rather
than by regenerating damaged descending
pathways~\cite{wernigMuller1992,dietz1998,behrmanHarkema2000}.

\paragraph{Which injury this model corresponds to.}
The brainstem-to-rhythm-generator (BS$\to$RG) drive stays tonic and non-zero
throughout, so the model structurally represents a state in which some
supraspinal descending input survives the lesion --- clinically,
motor-incomplete SCI (ASIA Impairment Scale grades C/D), the population in
which BWSTT and robot- or manually-assisted training reliably restore
stepping~\cite{wernigMuller1992,dietz1998,behrmanHarkema2000}. Two kinds of
synapse are plastic, and they map onto different mechanisms. CUT$\to$RG-E is
intraspinal and caudal to the lesion; its strengthening models
activity-dependent reinforcement of an intact, underused sensorimotor loop ---
the mechanism BWSTT and related therapies are thought to exploit. BS$\to$RG is
the model's spared descending projection; its strengthening is the analogue of
the reorganisation of spared reticulospinal input that accompanies spontaneous
locomotor recovery after incomplete SCI in rats~\cite{ballermannFouad2006}.
Neither models repair of a severed pathway.

An additional clinically relevant observation is that the effectiveness of
recovery depends on the degree of plasticity. In the simulations, too little plasticity leaves the gait on its safety timer
and an intermediate, still-developing level of plasticity produces the least
coordinated gait of all (\autoref{sec:results:comparison}); the fastest rate
converges within seconds but leaves the learned weights fluctuating from step
to step below their plateau, with weaker left--right alternation in fast walk.
The useful range of plasticity would thus correspond to rates that let spinal
motor networks complete their adaptation within the training period while
avoiding step-to-step instability, a heuristic that may help in choosing
rehabilitation intensities and in managing the effects of pharmacological or
electrical neuromodulation therapies on spinal adaptation.

\paragraph{Interpreting the learning-rate axis as plasticity capacity.}
One natural reading of the $\lambda$ axis is developmental. The capacity
for activity-dependent synaptic change is not fixed across the lifespan:
it is greatest during early postnatal critical periods and declines with
maturity, and correspondingly, neonatal and juvenile rodents recover
locomotor function after spinal injury considerably better than adults.
It is therefore tempting to read our fastest learning rates as
``younger'' and our slowest as ``older'' animals.\todo{cite developmental
critical-period / age-related decline in spinal plasticity and
age-dependent locomotor recovery after SCI} We suggest this analogy
should be treated as a qualitative ordering rather than a calibration:
$\lambda$ is a phenomenological STDP rate constant, it was not fitted to
any measurement of plasticity in animals of a given age, and age is only
one of several determinants of plasticity capacity --- time since injury,
prior activity history, and pharmacological or electrical neuromodulation
all modulate the same underlying quantity. Read this way, the $\lambda$
axis is better described as spanning \emph{plasticity capacity}, of which
age is one contributor.

With that caveat, the interpretation yields a concrete and testable
prediction. Because gait quality in our sweep does not increase
monotonically with $\lambda$ --- it is worst at an intermediate, incompletely
converged rate, and at the fastest rate the learned weights stay unstable
(\autoref{sec:results:comparison}) --- the model predicts that recovery
depends on plasticity being sufficient to \emph{complete} adaptation within
the training period, and that more plasticity is not automatically better. This resonates with the clinical picture in
which excessive or misdirected plasticity after SCI manifests as
maladaptive outcomes such as spasticity, aberrant sprouting and
neuropathic pain rather than improved locomotion.\todo{cite maladaptive
plasticity after SCI} The prediction is directly testable by running the
same rehabilitation protocol across age groups and asking whether
intermediate-age animals show the best coordination recovery, which would
be a stronger validation of the learning-rate axis than any of the
comparisons we are currently able to make.

Furthermore, the present computational framework might serve as a useful
tool for designing, optimising, and validating interventions before their
clinical application.

It is feasible to investigate, via simulation studies, various clinical
rehabilitation interventions (e.g., robot-assisted walking at various
assistance levels and weight-supported environments, ESCS at different
locations and frequencies, combination therapies) to identify the most
effective, non-invasive treatments. Finally, integrating this framework
with individual patient-specific data might pave the way for the
development of virtual-twin approaches that could predict, simulate, and
optimise patient outcomes after specific rehabilitation therapies in a
highly personalised fashion.

\subsection{Limitations}
In addition to the positive implications outlined above, it is important
to note some limitations of the proposed model. To simplify the biological
complexity of the spinal locomotor system, a level of abstraction was used
in defining sensorimotor pathways, interneurons and motor outputs. The
model does not include an adaptive supraspinal control mechanism (the
brainstem drive is a tonic input), detailed muscle-tendon
biomechanics beyond the saturating force/activation/fatigue model used here,
or longer-term structural plasticity.

It also omits features of SCI such as chronic damage to the dorsal and
ventral roots, glial scarring, or the diversity of human SCI.

It should also be highlighted that the model was not calibrated on any
experimental data collected from actual individuals with SCI. Hence the
values of the STDP parameters can only be viewed as computational
representations of plasticity rather than direct measures of spinal
learning capacity. Consequently, the present model should be interpreted
as hypothesis-generating rather than predictive in nature, and further
studies incorporating more detailed muscle dynamics, more realistic
sensory transduction, supraspinal inputs, and patient-specific data will
be required to establish greater biological realism.

These conceptual limitations are compounded by several specific modelling
choices worth stating explicitly:

\begin{enumerate}
  \item \textbf{Interneuron and interlimb detail.} We resolve two hand-set
        commissural projections rather than the genetically identified
        V0/V1/V2a/V3 interneuron classes, and we model one hindlimb pair
        rather than fore--hindlimb coordination~\cite{zhang2022}; we
        likewise do not reproduce gait-transition phenomena
        (walk $\to$ trot $\to$ gallop $\to$ bound).
  \item \textbf{Swing is timed, not sensed.} Only the end of stance is
        force-triggered; the model has no hip-position or other swing
        termination signal, so every swing lasts its timeout and
        stride-to-stride variability arises from stance alone.
  \item \textbf{Stride periods.} The measured strides ($\approx0.6$\,s fast
        walk to $\approx1.0$\,s slow walk) are long for a trotting rat,
        because every swing lasts its full timeout; locomotion modes are
        therefore defined by their operating points (Table~\ref{tab:modes}),
        not by treadmill speed.
  \item \textbf{Full unloading not reproduced.} At air-stepping loading the
        extensor force is too small for the force-triggered stance detector,
        so the loading axis stops at toe stepping.
  \item \textbf{Per-mode operating points.} Fatigue and timeout constants
        were chosen per mode, and a small persistent left--right fatigue
        asymmetry is needed at the fast and toe-stepping points to keep the
        legs from falling into in-phase stepping.
  \item \textbf{No stochastic replication.} Each condition is a single run
        and multi-threaded runs are not bit-reproducible; the ten-point
        initial-weight sweep (\autoref{sec:results:robustness}) substitutes
        for repeated stochastic trials in establishing that a result is not
        an artefact of a particular initialisation.
\end{enumerate}

\subsection{Concluding remark}
A computational sCPG half-centre oscillator, closed in the loop with
motoneurons, muscle proxies, and load-dependent sensory feedback that itself
decides when stance ends, tunes its own cutaneous and descending gains through
spike-timing-dependent plasticity to a stable set-point --- one that is
independent of how the network starts and of walking speed. Learning is what
turns the sensory loop on: without it the gait falls back on a safety timer,
and a partially learned synapse produces the least coordinated gait of all.
Partial unloading slows that learning and weakens the gait without abolishing
it. Beyond the specific findings, the model
demonstrates that a spinal CPG's descending and sensory gains need not be
prescribed by the modeller: given a plastic synapse, appropriately timed
sensory drive, and enough time, the network finds them on its own. That
property, together with the model's explicit and inspectable parameters,
is what makes it a plausible substrate for exploring --- and eventually,
with patient-specific calibration, for helping to design --- sensorimotor
rehabilitation strategies after spinal cord injury.
